\section{Trabajos Relacionados}

Esta sección presenta una revisión de trabajos previos relacionados con la clasificación de medicamentos mediante inteligencia artificial y sistemas de recomendación farmacéuticos.

\subsection{Clasificación de medicamentos mediante visión por computadora}

\subsubsection{Trabajo 1: Drug Identification System usando CNN}

Kumar et al. (2019) desarrollaron un sistema de identificación de medicamentos mediante deep learning que utiliza CNN para clasificar píldoras y tabletas a partir de imágenes.

\textbf{Tecnología utilizada:}
\begin{itemize}
	\item Red Neuronal Convolucional (CNN) con 6 capas convolucionales
	\item Transfer learning con ImageNet
	\item Framework: TensorFlow y Keras
\end{itemize}

\textbf{Resultados:}
\begin{itemize}
	\item Accuracy: 94.3\%
	\item Tiempo de inferencia: 0.8 segundos por imagen
\end{itemize}

\subsubsection{Trabajo 2: Pill Recognition usando ResNet-50}

Zeng et al. (2020) utilizaron ResNet-50 con fine-tuning específico para reconocimiento de medicamentos.

\textbf{Tecnología utilizada:}
\begin{itemize}
	\item ResNet-50 pre-entrenada
	\item Data augmentation avanzado
	\item Framework: PyTorch
\end{itemize}

\textbf{Resultados:}
\begin{itemize}
	\item Accuracy: 96.7\%
	\item Top-5 accuracy: 99.1\%
\end{itemize}

\subsection{Sistemas de recomendación farmacéuticos}

\subsubsection{Trabajo 3: Drug Recommendation usando Sentiment Analysis}

Doshi et al. (2021) propusieron un sistema de recomendación basado en análisis de sentimientos y componentes activos.

\textbf{Tecnología utilizada:}
\begin{itemize}
	\item Natural Language Processing (NLP)
	\item K-Nearest Neighbors para medicamentos similares
	\item Filtrado colaborativo
\end{itemize}

\textbf{Resultados:}
\begin{itemize}
	\item Precisión en recomendaciones: 87.5\%
	\item Satisfacción de usuarios: 78\%
\end{itemize}

\subsubsection{Trabajo 4: Content-Based Drug Recommendation}

Park et al. (2022) desarrollaron un sistema basado en similitud de componentes activos usando autoencoders.

\textbf{Tecnología utilizada:}
\begin{itemize}
	\item Autoencoder para embeddings de medicamentos
	\item Similitud coseno en espacio latente
	\item Base de datos FDA
\end{itemize}

\textbf{Resultados:}
\begin{itemize}
	\item Precisión: 91.2\%
	\item Recall: 85.7\%
\end{itemize}

\subsection{Diferenciación del proyecto MediFinder}

MediFinder integra ambas líneas de investigación:

\begin{itemize}
	\item Pipeline completo: desde captura de imagen hasta recomendación
	\item CNN para clasificación + KNN para recomendación
	\item Contextualizado a Colombia
\end{itemize}

\pagebreak